\documentclass[10pt]{article}
\usepackage[letterpaper,margin=0.72in]{geometry}
\usepackage[T1]{fontenc}
\usepackage[utf8]{inputenc}
\usepackage{amsmath,amssymb}
\usepackage{booktabs}
\usepackage{array}
\usepackage[hidelinks]{hyperref}

\hypersetup{
  pdftitle={Extended Supplementary Results and Reproducibility Details},
  pdfauthor={},
  pdfsubject={Supplementary reproducibility appendix},
  pdfkeywords={}
}

\setlength{\parindent}{0pt}
\setlength{\parskip}{5pt}
\pagestyle{plain}

\begin{document}

\begin{center}
{\Large\bfseries Extended Results and Reproducibility Details}\\
\vspace{2pt}
{\small Supplementary appendix for anonymous review}
\end{center}

\section*{Extended Results and Visualizations}

The supplementary package includes visualizations and extended tables tied to
analyses reported in the manuscript, including results that support the paper
but were omitted from the main text for space. Broader project figures,
exploratory outputs, notebooks, raw prediction caches, logs, and checkpoints
are excluded.

\begin{tabular}{>{\raggedright\arraybackslash}p{0.32\linewidth}>{\raggedright\arraybackslash}p{0.58\linewidth}}
\toprule
\textbf{Folder or file group} & \textbf{Contents and purpose} \\
\midrule
\texttt{extended\_visualizations/} & Figure PDFs directly tied to reported experiments and claims: OOD performance, blank-text ablation, counterfactual flip rate, mean probability shift, CKA, qualitative case studies, cross-attention visualization, DDI fairness, Holm-corrected shuffle significance, biopsy examples, SAE features, causal patching, matched grounding plots, and LVLM diagnostics. \\
\texttt{extended\_tables/} & Compact CSV/JSON outputs for reported manuscript claims: PAD-UFES OOD metrics, MILK10k lesion-disjoint audit and split-distribution check, metadata-shuffle control, DDI fairness, activation patching, SAE feature discovery, Feature 1449 knockout, DDI demographic patching, LVLM steering diagnostics, center-crop control, and matched grounding audit. \\
\bottomrule
\end{tabular}

\section*{Implementation and Reproducibility Details}

\subsection*{Datasets and Splits}

MILK10k is used as the in-domain training benchmark with six classes:
MEL, BCC, SCC, ACK, NEV, and SEK. The included \texttt{milk10k\_train.csv}
contains 9,344 image rows corresponding to 4,672 lesion identifiers, with
two images per lesion. The deterministic clean split groups rows by lesion
identifier, sorts lesion IDs, and reserves the final 15\% of lesions for
validation. This yields 700 validation lesions, 1,400 validation images,
7,944 training images, and zero lesion overlap. This split is used as a
leakage-control audit; repeated stratified lesion-disjoint splits would be the
appropriate robustness extension for in-domain performance estimation.
As an additional distribution check, we sampled 10,000 random lesion-disjoint
validation sets of the same size. The deterministic validation-set class
proportions were close to random expectation for every class (maximum absolute
$z$-score $0.71$), indicating that the fixed sorted-ID split is not an obvious
class-composition outlier. The CSV
\path{extended_tables/task29_lesion_split_distribution_audit.csv}
contains the per-class counts, random 95\% intervals, and tail probabilities.

We also repeated the Text-Only leakage check across five random stratified
lesion-disjoint splits using cached frozen BioClinicalBERT CLS embeddings and
the same linear-head training protocol. Text-Only accuracy averaged 46.28\%
(95\% CI 33.23--59.32), below the validation majority baseline of 53.92\%;
the mean Text-Only-minus-majority difference was $-7.65$ percentage points
($p=0.179$), providing no evidence that text-only generalization persists after
lesion overlap is removed. The per-split CSV and summary JSON are included as
\path{extended_tables/task30_textonly_random_lesion_cv.csv} and
\path{result_summaries/task30_textonly_random_lesion_cv_summary.json}.

PAD-UFES-20 is used for out-of-distribution evaluation through
\texttt{pad\_ufes\_20\_test.csv}. DDI is used for Fitzpatrick skin-tone
stratified analysis; raw DDI images are not included, but the audit code and
compact output tables are included. Raw public datasets and trained model
weights are intentionally excluded from the zip.

\subsection*{Preprocessing}

Images are resized to $224\times224$ and normalized with ImageNet statistics:
mean $(0.485, 0.456, 0.406)$ and standard deviation
$(0.229, 0.224, 0.225)$. Training augmentation uses random resized crop
(scale 0.8--1.0, aspect ratio 0.9--1.1), horizontal and vertical flips
with probability 0.5, random rotation up to 45 degrees, and color jitter
with brightness, contrast, and saturation 0.1 and hue 0.05. Evaluation uses
deterministic resize, tensor conversion, and normalization.

Clinical text is tokenized with \texttt{Bio\_ClinicalBERT} using maximum
length 128. Missing or literal null text is mapped to the configured blank
string. Labels use the six-class mapping MEL=0, BCC=1, SCC=2, ACK=3,
NEV=4, SEK=5.

\subsection*{Training Hyperparameters}

\begin{tabular}{ll}
\toprule
\textbf{Setting} & \textbf{Value} \\
\midrule
Optimizer & AdamW \\
Loss & Cross-entropy \\
Learning rate & $2\times10^{-5}$ \\
Weight decay & 0.01 \\
Batch size & 16 \\
Epochs & 5 maximum \\
Early stopping & validation loss patience of 2 epochs \\
Warmup & linear schedule with 10\% warmup ratio \\
Gradient clipping & maximum norm 1.0 \\
Mixed precision & enabled on CUDA \\
Seeds & 456, 789, 1337 for the main reported multi-seed runs \\
\bottomrule
\end{tabular}

\subsection*{Architecture Specifications}

All discriminative models use frozen ViT-Base-patch16-224 visual features
and frozen Bio\_ClinicalBERT text features, both with 768-dimensional hidden
states. The output space has six diagnostic classes.

\begin{tabular}{>{\raggedright\arraybackslash}p{0.28\linewidth}>{\raggedright\arraybackslash}p{0.62\linewidth}}
\toprule
\textbf{Model} & \textbf{Specification} \\
\midrule
Image-Only & Frozen ViT-B/16 encoder followed by a linear six-class classifier. \\
Text-Only & Frozen Bio\_ClinicalBERT CLS embedding followed by a linear six-class classifier. \\
Late Fusion & Concatenates 768-d visual and 768-d text embeddings, applies a 512-d hidden layer, ReLU, dropout 0.1, and a six-class classifier. \\
GMU & Projects visual and text embeddings to 512-d, computes a 512-d sigmoid gate from the concatenated 1536-d input, and classifies the gated fused representation. \\
Cross-Attention V-to-T & Uses one multi-head attention layer with 8 heads and dropout 0.1 where a visual query attends to text keys and values. \\
Cross-Attention T-to-V & Uses one multi-head attention layer with 8 heads and dropout 0.1 where text queries attend to ViT patch-token keys and values; the attended sequence is mean-pooled before classification. \\
Top-K SAE & Trained on the mean-pooled fused residual stream of Cross-Attention T-to-V with hidden dimension 6144 and $k=32$. \\
\bottomrule
\end{tabular}

\section*{Mathematical Derivations}

\subsection*{Lesion Leakage Under Image-Level Random Splitting}

MILK10k contains two images per lesion. Under an image-level random split
with validation probability $p_{\mathrm{val}}=0.15$, conditioning on one image
from a lesion appearing in validation, the probability that the paired image
appears in training is
\[
P(\text{paired image in train}\mid\text{one image in validation})
=1-p_{\mathrm{val}}=0.85.
\]
Thus an image-level split is expected to leak approximately 85\% of validation
lesions into training. Grouping by lesion ID before splitting removes this
failure mode by construction:
\[
\mathcal{L}_{\mathrm{train}}\cap\mathcal{L}_{\mathrm{val}}=\varnothing.
\]

\subsection*{Counterfactual Flip Rate}

For image $x_i$, real text $t_i$, and contradictory text $\tilde{t}_i$, let
\[
\hat{y}_i = \arg\max_c f_c(x_i,t_i),
\qquad
\tilde{y}_i = \arg\max_c f_c(x_i,\tilde{t}_i).
\]
The counterfactual flip rate is
\[
\mathrm{CFR}=\frac{1}{N}\sum_{i=1}^N
\mathbf{1}\{\hat{y}_i\neq \tilde{y}_i\}.
\]
For an image-only model $f(x,t)=g(x)$, predictions do not depend on $t$.
Therefore $\hat{y}_i=\tilde{y}_i$ for all $i$, and
\[
\mathrm{CFR}_{\mathrm{image-only}}=0.
\]
This is why CFR against an image-only baseline is a structural property rather
than evidence of semantic grounding.

\subsection*{Linear CKA Identity for Identical Feature Spaces}

Linear CKA between centered feature matrices $X,Y\in\mathbb{R}^{N\times d}$ is
\[
\mathrm{CKA}(X,Y)=
\frac{\|X^\top Y\|_F^2}{\|X^\top X\|_F\|Y^\top Y\|_F}.
\]
If the fused representation is identical to the visual representation, then
$Y=X$ and
\[
\mathrm{CKA}(X,X)=
\frac{\|X^\top X\|_F^2}{\|X^\top X\|_F\|X^\top X\|_F}=1.
\]
Thus CKA equal to 1.0 is guaranteed for an unchanged image-only feature space
and must not be interpreted as a semantic text-grounding result.

\subsection*{Matched Grounding Margin}

For a binary class pair such as MEL/NEV, define the class margin under text
condition $t$ as
\[
m_i(t)=f_{\mathrm{MEL}}(x_i,t)-f_{\mathrm{NEV}}(x_i,t).
\]
For each matched image, the grounding shift is
\[
\Delta_i=m_i(t_i^{\mathrm{aligned}})-m_i(t_i^{\mathrm{contradictory}}).
\]
Positive $\Delta_i$ means aligned text increases the true-pair margin relative
to contradictory text. The reported real-image mean shift is negative for the
representative MEL/NEV audit, indicating that aligned text did not improve the
correct-class margin.

\subsection*{SAE Feature Knockout Estimand}

Let $z_i$ be the fused representation, $a_{ij}$ the activation of SAE feature
$j$, and $d_j$ its decoder direction. A single-feature knockout constructs
\[
z_i^{(-j)}=z_i-a_{ij}d_j.
\]
The reported intervention estimand for Feature 1449 is the mean change in
malignant-minus-benign logit contrast,
\[
\Delta_{\mathrm{mal-ben}}=
\frac{1}{N}\sum_{i=1}^N
\left[
\left(\overline{f}_{\mathrm{mal}}(z_i^{(-1449)})
-\overline{f}_{\mathrm{ben}}(z_i^{(-1449)})\right)
-
\left(\overline{f}_{\mathrm{mal}}(z_i)
-\overline{f}_{\mathrm{ben}}(z_i)\right)
\right].
\]
The random active-feature control uses the same estimator with a non-target
active feature selected per row. Because decoded SAE-direction subtraction can
move representations off manifold and can have non-local downstream effects,
this estimand should be interpreted as causal-intervention evidence for a
candidate shortcut feature, not as proof that a single isolated clinical artifact
fully explains the classifier's behavior.

\subsection*{Loss Function Note}

The discriminative training objective is standard multi-class cross-entropy:
\[
\mathcal{L}=-\frac{1}{N}\sum_{i=1}^N \log
\frac{\exp f_{y_i}(x_i,t_i)}{\sum_{c=1}^{6}\exp f_c(x_i,t_i)}.
\]
No novel loss function is introduced in the manuscript; the mathematical
contributions are the audit definitions, controls, and derivations above.

\end{document}
